\documentclass[12pt]{article}
\usepackage[utf8]{inputenc}
\usepackage[T1]{fontenc}
\usepackage[brazilian]{babel}
\usepackage{amsmath}
\usepackage{geometry}
\geometry{margin=2.5cm}

\title{Exemplo-Modelo --- Questão 01, Capítulo 1}
\author{ECOM028 --- Repositório de Soluções (arquivo de referência)}
\date{}

\begin{document}
\maketitle

\section*{Enunciado}

Considere uma pilha real com tensão em circuito aberto
$V_{oc} = 1{,}56\,\text{V}$. A resistência de carga é calculada a partir dos
dois últimos dígitos da matrícula do aluno:
\[
R_{carga}(\Omega) = \big(\text{2 últimos dígitos da matrícula} \bmod 9\big) + 1.
\]

Para este exemplo-modelo, adota-se uma matrícula fictícia terminada em
\textbf{58}, o que fornece:
\[
58 \bmod 9 = 4 \quad \Rightarrow \quad R_{carga} = 4 + 1 = 5\,\Omega.
\]

\textbf{Pergunta:} determine $P_{alvo}$, definida como a maior potência (em
mW) para a qual
\[
V_{op} = \sqrt{P_{alvo} \times R_{carga}}
\]
cai entre $0{,}90\,\text{V}$ e $1{,}56\,\text{V}$, utilizando as curvas de
descarga em potência constante fornecidas no Capítulo 1 do e-book
(figura ``Coppertop AA --- Constant Power'').

\section*{Conceitos mobilizados por esta questão}

Ver arquivo \texttt{conceitos.yaml} nesta mesma pasta para o checklist
formal. Em resumo, esta questão trabalha: modelo de fonte real
($V_{oc}$, $R_{int}$), leitura de curvas de descarga em potência constante,
e o conceito de ponto de operação $(V_{op}, I_{op})$.

\end{document}
